% ---
% id: meat-core/rem:representability-ceiling
% kind: rem
% title: "The representability ceiling"
% summary: >-
%   The condition N ≤ 53 + log₂ϵ of [[prop:optol-nobound]] is what exact
%   admissibility forces on a two-outcome binary64 family. It caps N at 52,
%   attained at ϵ = 1/2, and belongs to the number model, not to the map from
%   the measurement to the rate.
% prominence: standard
% status: accepted
% uses: [def:declaration, prop:value-range]
% mentions: [prop:optol-nobound, rem:classes, req:infinite-value]
% assumes: []
% refs: []
% provenance: newly-derived
% lean:
% ---

\begin{remark}\label{rem:representability-ceiling}
Proposition~\ref{prop:optol-nobound} needs a declared family $(g_0,g_1)$
with $g_0+g_1=1$ and both entries binary64: exact admissibility is an
identity of the declared numbers (Definition~\ref{def:declaration}, the
exact-admissibility paragraph), not of the reals they approximate, and it
additionally requires the one addition of \eqref{eq:declorder} to be
exact. The smaller entry $g_0$ satisfies $g_0\le\tfrac12$, so
$g_1=1-g_0$ lies in the binade $[\tfrac12,1)$, the binary64 numbers with
exponent $-1$, where consecutive numbers are $2^{-53}$ apart (the unit in
the last place); hence $g_1$ is representable only if $g_0$ is a multiple
of $2^{-53}$, and $g_0>0$ makes that $g_0\ge2^{-53}$. With
$g_0=\epsilon2^{-N}$ this reads $\epsilon2^{-N}\ge2^{-53}$, that is
$2^{N}\le\epsilon\,2^{53}$, the inequality the proposition assumes of
$N$ and $\epsilon$; since $\epsilon\le\tfrac12$, it gives $N\le52$. The
value is attained: $\epsilon=\tfrac12$ and $N=52$ give
$(2^{-53},1-2^{-53})$, with the addition exact. The proof builds its
families at the largest power of two $\epsilon'\le\epsilon$, with
$g_0=\epsilon'2^{-N}$, and loses no $N$ by it: $2^{N-53}$ is itself a
power of two, so the inequality holds at $\epsilon'$ exactly when it
holds at $\epsilon$.

The ceiling is a property of the number model at these fixed dimensions,
not of the map $\{F_v\}\mapsto\hlow$: in exact real arithmetic the same
construction runs at every $B>0$, and it is the number model that stops
it. What the number model does is cap how much of the gap a family of
finite value, $g_0>0$, can express, at $\frac{\alpha}{\alpha-1}N$ with
$2^{N}\le\epsilon\,2^{53}$; the exactly admissible family $(0,1)$ has the
value $\hlow=+\infty$ (Proposition~\ref{prop:value-range}\,(iv)) and carries
no key-length claim (Requirement~\ref{req:infinite-value}), so it is set
aside. At $\epsilon=2^{-40}$ the instance still
delivers $N=13$, thirteen bits per round or more, and it runs down to
$\epsilon=2^{-52}$, where $N=1$ still gives $\frac{\alpha}{\alpha-1}>1$
bit per round. At the instance's $d_Z=1$ both exceed
$\log_2|\Zset|=1$ bit per round, the entropy-part ceiling of
Proposition~\ref{prop:value-range}\,(ii): more than the entropy terms of $\hlow$ can
contribute at any declaration with $d_Z=1$, the rest of $\hlow$ being the
acceptance-test penalty. The comparison is the instance's: at a larger
$d_Z$ the entropy-part ceiling $\log_2(d_Z+1)$ can exceed a gap of that
size (at $\alpha=1.9$, $d_Z=7$ and $N=1$ it is $3$, against
$\frac{\alpha}{\alpha-1}N\approx2.1$). The representability ceiling
shapes part~(ii) of the proposition, which rules out allowances below the
exhibited scale rather than every allowance. It shapes as well
Remark~\ref{rem:classes},
which cites part~(i) to explain why a tolerance may refuse the deletion of
a sector but never license one, and which states the gap there as
$\frac{\alpha}{\alpha-1}N$ with $N$ up to $52$. The conclusion that
matters is what the bounded statement still delivers: on the data of
part~(i), part~(ii) with $N=1$ refutes, at every
$\epsilon\in[2^{-52},\tfrac12]$, every $\mathrm{err}$ below
$\frac{\alpha}{\alpha-1}$, so every $\mathrm{err}$ at or below that
entropy-part ceiling of $1$ bit per round. The cap moves the refuted range
from every allowance to the allowances below
$\frac{\alpha}{\alpha-1}\cdot52$, still far above that ceiling, and no
rule of this module turns on the difference.
\end{remark}
